\documentclass{article}
\usepackage{graphicx} % Required for inserting images
\usepackage{amsmath}
\usepackage{float}
\title{Lesson 9 Homework - Calculus}
\author{William Wang}
\date{September 2026}

\begin{document}

\maketitle

\section{}

\subsection{}

The slope of the line through two points is the change in $y$ over the change in $x$:
\[
m_{PQ} = \frac{f(b)-f(a)}{b-a}
\]

\subsection{}

Substituting $b = a+h$, we get
\[
m_{PQ} = \frac{f(a+h)-f(a)}{(a+h)-a} = \frac{f(a+h)-f(a)}{h}
\]

\subsection{}

Since $Q = (a+h, f(a+h))$ and $P = (a, f(a))$, the point $Q$ approaches $P$ exactly when its $x$-coordinate $a+h$ approaches $a$.
\subsection{}
\[
f'(a) = \lim_{h \to 0} \frac{f(a+h)-f(a)}{h}
\]

For every fixed $h \neq 0$, the difference quotient $\frac{f(a+h)-f(a)}{h}$ is the slope of some secant line through $P$ and a nearby point $Q$. As $h \to 0$, the point $Q$ slides along the curve toward $P$, and the secant lines $PQ$ rotate, settling toward the line that best approximates the curve's direction right at $P$,  that line is, by definition, the tangent line.

\subsection{}

Let 
\[f(x) = \frac{|h|}{h}\]
By definition:
\[
f'(1) = \lim_{h \to 0} \frac{f(1+h)-f(1)}{h} = \lim_{h \to 0} \frac{|1+h-1| - |1-1|}{h} = \lim_{h \to 0} \frac{|h|}{h}
\]

Now check the two one-sided limits, since $|h|/h$ behaves differently depending on the sign of $h$:

\begin{itemize}
    \item As $h \to 0^{+}$: $h>0$, so $|h| = h$, giving $\dfrac{|h|}{h} = \dfrac{h}{h} = 1$.
    \item As $h \to 0^{-}$,  $h<0$, so $|h| = -h$, giving $\dfrac{|h|}{h} = \dfrac{-h}{h} = -1$.
\end{itemize}

Since
\[
\lim_{h \to 0^{+}} \frac{|h|}{h} = 1 \neq -1 = \lim_{h \to 0^{-}} \frac{|h|}{h}
\]
the two one-sided limits disagree, so $\displaystyle\lim_{h\to 0} \frac{|h|}{h}$ doesn't exist. Therefore $f'(1)$ does not exist.
Geometric interpretation: this makes sense since the graph of $|x-1|$ has a sharp corner at $x=1$. Approaching from the right, secant lines have slope trending to $+1$; approaching from the left, they trend to $-1$. Since the secant slopes approach different values depending on direction, there's no single well-defined tangent line at the corner which is consistent with the limit failing to exist.

\section{}
\begin{figure}[H]
    \centering
    \includegraphics[width=1\linewidth]{image.png}
\end{figure}
The graph appears to have a vertical tangent line at \(x = 0\). To confirm this, we take the limit as \(h\to 0\). We have to calculate the right hand and left hand derivative separately. 
\[\lim_{h\to 0^+} \frac{f(0+h) - f(0)}{h} = \lim_{h \to 0^+} \frac{\sqrt{h}}{h} = \lim_{h\to 0^+} \frac{1}{\sqrt{h}}=  \infty\]
\[\lim_{h\to 0^-} \frac{f(0+h) - f(0)}{h} = \lim_{h \to 0^-} \frac{-\sqrt{-h}}{h} = \lim_{h\to 0^-} = \frac{1}{\sqrt{-h}} = \infty\]
Since the left hand and right hand derivatives agree with each other, the limit is \(\infty\). This confirms a vertical tangent line at \(x= 0 \)
\section{}
We have the function 
\[f(x) = \begin{cases}
    x^2 \ \text{if } x< 0 \\
    x  \ \text{if } x \geq 0
\end{cases}\]
Then the right hand derivative is 
\[\lim_{h \to 0^+} \frac{f(0+h) - f(0)}{h} = \lim_{h\to 0^+}\frac{h}{h} = 1\]
The left hand derivative is 
\[\lim_{h\to 0^-} \frac{f(0+h) - f(0)}{h} = \lim_{h\to 0^-} \frac{h^2}{h} = \lim_{h\to 0^-} h = 0\]
Because the left hand derivative is not equal to the right hand derivative, it is not differentiable at \((0,0)\)



\end{document}
